\documentclass{article}
\usepackage{amsmath,amssymb}
\usepackage{xurl}
\usepackage[hidelinks]{hyperref}
% Shared commands for notes in docs/paper-gaps/.

\DeclareUnicodeCharacter{2080}{\ifmmode{}_{0}\else\textsubscript{0}\fi}
\DeclareUnicodeCharacter{2081}{\ifmmode{}_{1}\else\textsubscript{1}\fi}
\DeclareUnicodeCharacter{2082}{\ifmmode{}_{2}\else\textsubscript{2}\fi}
\DeclareUnicodeCharacter{2099}{\ifmmode{}_{n}\else\textsubscript{n}\fi}

\newcommand{\Lean}{\textsc{Lean}}
\ExplSyntaxOn
\NewDocumentCommand{\leanid}{m}{
  \ifmmode
    \textsf{\detokenize{#1}}
  \else
    \group_begin:
    \urlstyle{sf}
    \tl_set:Nn \l_tmpa_tl {#1}
    \tl_replace_all:Nnn \l_tmpa_tl {\_} {_}
    \exp_args:NV \path \l_tmpa_tl
    \group_end:
  \fi
}
\ExplSyntaxOff
\providecommand{\tr}{\operatorname{tr}}
\newcommand{\tnleanIssueBase}{https://github.com/LionSR/TNLean/issues/}
\newcommand{\tnleanPrBase}{https://github.com/LionSR/TNLean/pull/}
\newcommand{\tnleanDocBase}{https://sirui-lu.com/TNLean/docs/}
\newcommand{\ghissue}[1]{\href{\tnleanIssueBase#1}{GitHub issue \##1}}
\newcommand{\ghpr}[1]{\href{\tnleanPrBase#1}{PR \##1}}
\newcommand{\doclink}[1]{\href{\tnleanDocBase#1.html}{\path{#1}}}

% Dirac notation and the identity operator, as in blueprint/src/macros/common.tex.
% \providecommand keeps note-local definitions authoritative.
\usepackage{dsfont}
\providecommand{\ket}[1]{|#1\rangle}
\providecommand{\bra}[1]{\langle #1|}
\providecommand{\braket}[2]{\langle #1 | #2 \rangle}
\providecommand{\ketbra}[2]{|#1\rangle\!\langle #2|}
\providecommand{\Id}{\mathds{1}}
\providecommand{\one}{\mathds{1}}
\providecommand{\C}{\mathbb{C}}
\providecommand{\MN}[1]{M_{#1}(\C)}
\providecommand{\MD}{M_D(\C)}

% Verdict marker: \gapnote{<kind>}{<status>}. Kind is the policy
% classification (clarification, local-correction, scope-restriction,
% unfaithful, false-source, open-gap); status is open, wip (elimination
% underway), resolved, or historical. Machine-read by the site index;
% prints a verdict line.
\newcommand{\gapnote}[2]{\par\noindent\textbf{Verdict:} #1 (#2).\par}

\title{The Domain of the Varying-Bond Open Parent Theorem}
\author{}
\date{October 29, 2026}
\begin{document}
\maketitle
\gapnote{scope-restriction}{open}

\section{Source and actual conclusion}
PGVWC07, arXiv:quant-ph/0608197, Theorem 9 and proof,
\path{Papers/quant-ph_0608197/MPSarchive.tex}, lines 1195--1237,
starts from a canonical open-boundary representation whose blocked site
families satisfy C1. The local interaction is the orthogonal projector
onto the complement of the corresponding two-site boundary range.

The current candidate constructs these actual projectors and their open
sum, derives the adjacent coefficient sections from the actual kernel,
reconstructs the full chain by the printed canonical adjoint method,
and closes the singleton endpoint bonds. For a common physical alphabet
and $N\ge2$, the conclusion is
\[
    H\ge0,\qquad \|\Omega\|=1,\qquad
    \ker H=\mathbb C\Omega,\qquad \dim\ker H=1.
\]
The unit norm and both kernel inclusions are derived from the supplied
site matrices. There is no assumed ground state, local annihilation
condition, larger-block boundary, or canonical-form supplier.

\section{Scope still retained}
The proof is stated for at least two sites, the domain with an adjacent
pair. It does not assert an unrestricted conclusion for lengths without
such a pair. The physical alphabet is common to the sites, as in the
existing finite open-chain data. A conclusion for differently sized
blocked physical spaces has not been claimed by this statement.
These restrictions must remain visible in its mathematical entry. The
unrestricted source theorem is not marked complete by this candidate.

The remaining scope work is to establish the intended short-length
convention from the full source setting and, if required, extend the
finite-word induction and coordinate restrictions to site-dependent
physical alphabets. No counterexample to the unrestricted source claim
is asserted here. The current proofs also await hosted Lean validation.

\section{Formalization}
The actual local terms and their kernel characterization are in
\path{TNLean/MPS/ParentHamiltonian/VaryingOBCProjectorKernel.lean}.
The full-chain adjoint induction is in
\path{TNLean/MPS/ParentHamiltonian/VaryingOBCReconstruction.lean},
and actual boundary-state annihilation is in
\path{TNLean/MPS/ParentHamiltonian/VaryingOBCBoundaryStates.lean}.
The scalar-endpoint normalization and final kernel statements are in
\path{TNLean/MPS/ParentHamiltonian/VaryingOBCParentUniqueness.lean}.

\end{document}
